% ======================================================================
%  PART VII İÇERİĞİ (devam): Yapı Taşları başvuru tablosu
% ======================================================================

\concept{Reference · Başvuru}{Yapı Taşları:\\Hangisi Ne Zaman?}{Tool · Skill · MCP · Agent · Subagent --- en sık karıştırılan dörtlü, net ayrımıyla}

En sık karışan kavramlar bunlardır --- çünkü \emph{iki farklı eksende} yaşarlar.
Doğru soru ``hangisi?'' değil, önce \textbf{hangi eksen?}:

\begin{itemize}
  \item \textbf{Yetenek verme} (model \emph{karar vermez}; kim çalıştırır / nasıl paketlenir / nasıl bağlanır): \term{Tool} · \term{Skill} · \term{MCP}
  \item \textbf{Karar devretme} (kararı \emph{modele bırakırsın}; döngü / delegasyon): \term{Agent} · \term{Subagent}
\end{itemize}

\vspace{2pt}
\noindent\begin{tabularx}{\linewidth}{@{}>{\raggedright\arraybackslash}p{24mm} L L@{}}
\arrayrulecolor{hairline}\toprule
{\sffamily\bfseries\small Yapı taşı} & {\sffamily\bfseries\small Ne} & {\sffamily\bfseries\small Ne zaman} \\
\midrule
\multicolumn{3}{@{}l}{\cellcolor{wash}{\sffamily\bfseries\footnotesize\color{accentdk}~YETENEK VERME · karar VERMEZ}} \\
\addlinespace[3pt]
\textbf{Tool} & Tek fonksiyon çağrısı & Canlı veri / kesin işlem \\
\arrayrulecolor{hairline}\cmidrule[0.2pt]{1-3}
\textbf{Skill} & Paketli yetkinlik (talimat + script) & Tekrar eden, paketlenebilir uzmanlık \\
\cmidrule[0.2pt]{1-3}
\textbf{MCP} & Araç bağlama standardı (protokol) & Araç çok uygulamada kullanılacak \\
\addlinespace[4pt]
\multicolumn{3}{@{}l}{\cellcolor{wash}{\sffamily\bfseries\footnotesize\color{accentdk}~KARAR DEVRETME · karar VERİR}} \\
\addlinespace[3pt]
\textbf{Agent} & Hedef + araç + döngü; adımları seçer & Adımlar önceden bilinemiyor \\
\cmidrule[0.2pt]{1-3}
\textbf{Subagent} & Alt-görevi devralan ayrı agent & Büyük görevi parçalara bölmek \\
\bottomrule
\end{tabularx}

\vspace{4pt}
{\sffamily\bfseries\small\color{accentdk}En sık karıştırılanlar:}
\begin{itemize}
  \item \textbf{MCP $\neq$ Agent} --- MCP bir bağlantı katmanıdır, karar vermez; bir \emph{workflow} da MCP kullanır. MCP'li olmak agent olmak demek değildir.
  \item \textbf{Skill $\neq$ Tool} --- Skill, talimat + kaynak içeren bir \emph{paket}tir; tool tek bir fonksiyondur. Bir skill içinde araçlar olabilir.
  \item \textbf{Skill $\neq$ Agent} --- Skill bir \emph{yetenektir}; agent onu \emph{kullanan döngüdür}.
  \item \textbf{Subagent $\neq$ Multi-Agent} --- Subagent bir parent'a bağlıdır; \term{Multi-Agent} eş düzey, koordine agent'lardır.
\end{itemize}

\begin{principlebox}
\textbf{Önce ekseni sor:} ``Kararı mı devrediyorum (agent · subagent), yoksa yetenek
mi veriyorum (tool · skill · mcp)?'' Çoğu karışıklık, bu iki ekseni birbirine
karıştırmaktan doğar --- ve ``her şeyi agent yapmak'' bunun en pahalı sonucudur.
\end{principlebox}
